\section{问题一：无噪声异步数据的时间对齐}

设两类定位数据均为同一条连续真实轨迹的无噪声离散采样。先分别构造连续轨迹插值，在公共时间域上以运动特征互相关给出多个粗候选，再用位置域均方误差
\begin{equation}
J(\delta)=\frac{1}{N(\delta)}\sum_{t_k\in\Omega(\delta)}
\left\|\bm r_1(t_k)-\bm r_2(t_k-\delta)\right\|_2^2
\end{equation}
进行全局粗扫与有界连续精修。互相关只用于缩小候选范围，最终时间偏差由位置 MSE 决定，从而降低重复轨迹形状产生伪相关峰的风险。

正式结果为 $\delta=\QOneDelta\,\mathrm{s}$，对齐 RMSE 为 $\QOneRMSE\,\mathrm{m}$，达到数值误差量级。对齐后将两类无噪声采样点合并，并在不跨异常大时间空洞的前提下重建 10 Hz 轨迹。

问题一额外通过多插值器比较、已知时间偏差合成试验、搜索范围与步长敏感性分析验证结果稳定性。由于该问题不存在随机测量噪声，近零对齐残差是模型正确性的直接检验之一。
